% ==========================================================================
% GENERATED by scripts/split_paper.py. Do not edit this file.
%
% Source: imported_papers/grokking/iclr2027_conference_v2.tex
% Part 8: The grain sweep and the test on published optimizers
%
% Edits belong upstream, in the published paper. To change what the
% chapter includes, change imports/grokking.toml and re-run
% the splitter.
% ==========================================================================
\section{The grain sweep and the test on published optimizers}

\subsection{The grain sweep}
\label{app:grain}
The credit knob equalizes credit one unit at a time, and the sweep applies the same operation to \emph{groups} of units, which is what a normalizer of any coarser grain effectively does. We split the $n = 256$ units of the first hidden layer into $g$ equal groups of $m = n/g$ units, and at every step rescale each group's gradient rows by the single factor that brings its total credit to the average group's, leaving the split \emph{within} the group untouched. Sweeping $m$ from $256$ down to $1$ carries the rule from equalizing nothing to equalizing every unit separately, and Proposition~\ref{prop:grain} derives the participation each $m$ should produce.

Both ends check themselves. At $m = 256$ the rescaling is the identity, so this end must reproduce backpropagation, and it does, at $0.602\,(0.023)$ against the $0.60\,(0.02)$ measured elsewhere in the paper. At $m = 1$ it is exactly Equation~\ref{eq:knob} at $\beta = 1$, so this end must reproduce the knob, and it does, at $0.992$ against $0.99$. Only the grains between them carry new information.

The steep end carries most of the effect. Half of the total movement arrives in the last halving, from $0.753$ at two units per group to $0.992$ at one, and targeted deletion tolerance rises along the same curve, from $0.327$ to $0.576$. Groups of two are still groups: a rule that comes within a factor of two of the right grain gets almost none of the effect.

The same shape holds on $S_5$, flat through $0.749$ at four units per group and rising to $0.852$ and $0.974$, so the result stands on a non-abelian task where it cannot be an artifact of the Fourier structure of modular arithmetic. Two points inside the flat region sit below their coarser neighbour, by less than one standard deviation in both cases, so we read the region as flat rather than as non-monotone.

\subsection{The prediction on published optimizers}
\label{app:zoo}
Proposition~\ref{prop:grain} says that equalizing credit moves the circuit only when it is done one unit at a time, so knowing the level at which an optimizer normalizes should be enough to say where its circuit lands, without knowing anything else about how it works. We test that on eight published optimizers spanning five grains, per layer, per matrix, factored, per parameter and per neuron, and all eight land where their grain says (Table~\ref{tab:zoo}).

\begin{table}[!ht]
\caption{The grain prediction holds for all eight published rules, $8/8$ in the \emph{branch} column. Published rules chosen by grain (three seeds): LARS \citep{you2017lars}, RMSProp \citep{tieleman2012rmsprop}, LAMB \citep{you2020lamb}, Adafactor \citep{shazeer2018adafactor}, Lion \citep{chen2023lion}, Muon \citep{jordan2024muon}, Shampoo \citep{gupta2018shampoo} and Nero \citep{liu2021nero}. \emph{Branch} is what Proposition~\ref{prop:grain} predicts from the rule's grain alone, and $\checkmark$ marks it satisfied. The two anchor rows are the ends the sweep must reproduce by construction, so they carry no prediction. Seven rules grok here, and Nero reaches test $0.824$ at this learning rate and groks at its tuned rates, landing distributed in both runs (Table~\ref{tab:named}). Every rule of a grain other than per unit stays in the concentrated range whatever its credit participation does, and the per-unit rules land distributed.}
\label{tab:zoo}
\begin{center}
\footnotesize
\setlength{\tabcolsep}{4pt}
\begin{tabular}{llcccl}
\toprule
Rule & Grain & Test & Credit PR & Circuit PR & Branch \\
\midrule
Backpropagation & --- (anchor)  & $0.998$ & $0.586$ & $0.595$ & $0.60$ (anchor) \\
LARS            & per layer     & $0.967$ & $0.520$ & $0.613$ & plain flow $\checkmark$ \\
RMSProp         & per parameter & $0.992$ & $0.791$ & $0.470$ & plain flow $\checkmark$ \\
LAMB            & layer $\times$ param & $0.995$ & $0.779$ & $0.543$ & plain flow $\checkmark$ \\
Adafactor       & factored      & $0.991$ & $0.839$ & $0.518$ & plain flow $\checkmark$ \\
Lion            & per parameter & $0.985$ & $1.000$ & $0.535$ & plain flow $\checkmark$ \\
Muon            & per matrix    & $0.999$ & $0.478$ & $0.524$ & plain flow $\checkmark$ \\
Shampoo         & per matrix    & $0.986$ & $0.521$ & $0.555$ & plain flow $\checkmark$ \\
\midrule
Nero            & per neuron    & $0.824$ & $1.000$ & $0.960$ & tracks credit $\checkmark$ \\
Knob ($\beta{=}1$) & per unit   & $0.976$ & $1.000$ & $0.992$ & tracks credit (anchor) \\
\bottomrule
\end{tabular}
\end{center}
\end{table}

Optimization geometry is known to select among equally accurate interpolators, from steepest and mirror descent under different norms \citep{gunasekar2018implicit} to the $\ell_1$ bias of small initialization \citep{woodworth2020kernel}. Those results say \emph{that} the update selects. Ours says \emph{which statistic} does the selecting, the per-unit credit profile, and that its grain is decisive, since normalized methods grouping per layer or per parameter \citep{you2017lars,you2020lamb,bernstein2018signsgd} leave the circuit concentrated. This is complementary to rich versus lazy \citep{kumar2024grokking}, which varies the parameterization scale rather than how a fixed rule distributes credit.
